% ===================================================================
%  Αρχείο: 1749_SOLUTION.tex (Fragment)
%  Αυτόματη μετατροπή μέσω doc2latex.py
% ===================================================================

ΘΕΜΑ 4

Θεωρούμε δυο σημεία Α και Β τα οποία βρίσκονται στο ίδιο μέρος ως προς μια ευθεία (ε), τέτοια ώστε η ευθεία ΑΒ δεν είναι κάθετη στην (ε). Έστω Α΄ το συμμετρικό του Α ως προς την ευθεία (ε), δηλαδή η (ε) είναι μεσοκάθετος του ΑΑ΄. $\ $

\begin{alist}
\item Αν η Α΄Β τέμνει την ευθεία (ε) στο σημείο Ο, να αποδείξετε ότι:

.}
\end{alist}
\begin{alist}
\item
  Η ευθεία (ε) διχοτομεί τη γωνία $\widehat{ΑΟΑ΄}$. \monades{6}
\item
  Οι ημιευθείες ΟΑ και ΟΒ σχηματίζουν ίσες οξείες γωνίες με την ευθεία (ε).

\monades{6}

\item Αν Κ είναι ένα άλλο σημείο πάνω στην ευθεία (ε), να αποδείξετε ότι:

.}
\end{alist}
\begin{alist}
\item
  ΚΑ=ΚΑ΄ \monades{6}
\item
  ΚΑ+ΚΒ\textgreater ΑΟ+ΟΒ \monades{7}
\end{alist
